\chapter{Formal Verification} \label{ch:verification}

This chapter describes the Alloy bounded model checking methodology used to verify
structural properties of the five covenant vault designs. Section~\ref{se:alloy_bg}
introduces Alloy and its applicability to protocol verification.
Section~\ref{se:properties_verified} defines the properties checked and reports
results. Section~\ref{se:three_layer} describes how the formal, empirical, and
economic evidence layers reinforce each other. The full model architecture and
per-covenant detailed results appear in Appendix~\ref{app:alloy_details}.

% ======================================================================
\section{Alloy and Bounded Model Checking} \label{se:alloy_bg}

Alloy~\cite{Jackson12} is a declarative modeling language with automatic bounded
analysis. An Alloy model specifies a universe of \emph{signatures} (types),
\emph{relations} (fields), and \emph{constraints} (facts and predicates). The Alloy
Analyzer exhaustively searches all instances within a user-specified \emph{scope}
(maximum number of atoms per signature) to find counterexamples to assertions or
satisfying instances for predicates.

Alloy's bounded model checking is well-suited to protocol verification for two reasons.
First, the small-scope hypothesis---if a property has a counterexample, it typically
has one within a small bound---means that modest scopes (6--10 atoms) provide high
confidence without exploring astronomically large state spaces. Second, Alloy's
relational logic naturally expresses state machines, ownership relations, and
authorization policies that characterize vault protocols.

dgpv~\cite{dgpv24} applied Alloy to verify a single OP\_CAT vault prototype
(Rijndael's design~\cite{Rij24}), checking fund conservation and covenant enforcement
within that one design. Our work extends this approach to all four BIP-specified
covenant protocols plus Simplicity, using a shared base model to ensure uniform
property definitions across designs. The models comprise 9 Alloy files: two base
modules (\texttt{btc\_base}, \texttt{vault\_base}), five concrete vault models,
a threat model module, and a cross-covenant interaction module.

% ======================================================================
\section{Properties and Results} \label{se:properties_verified}

We verify three categories of structural properties across the five vault designs.

\paragraph{Fund conservation.}
No transaction in the vault lifecycle creates value:

\begin{equation} \label{eq:fund_conservation}
\forall\, t \in \text{Transactions}:
  \sum_{\text{out} \in t.\text{outputs}} \text{out.amount}
  \leq \sum_{\text{in} \in t.\text{inputs}} \text{in.amount}
\end{equation}

\paragraph{No-extraction.}
An attacker with a specified key set cannot increase their holdings through any
sequence of vault transitions:

\begin{equation} \label{eq:no_extraction}
\forall\, \text{path} \in \text{Executions},\;
\forall\, a \in \text{Attackers}:
  a.\text{holdings}_{\text{final}} \leq a.\text{holdings}_{\text{initial}}
\end{equation}

\noindent
We check this under three attacker profiles per design: no keys, hot key only, and
cold key only. The one expected failure is CAT+CSFS under cold-key-only compromise:
the Alloy model finds a counterexample where the cold key holder redirects recovery
funds to an arbitrary address, structurally confirming TM11
(Section~\ref{se:cold_key_recovery}).

\paragraph{Recovery liveness.}
From any reachable vault state, a recovery transition exists:

\begin{equation} \label{eq:recovery_liveness}
\forall\, s \in \{\text{Funded}, \text{Triggered}\}:
  \exists\, t \in \text{Transitions}:\;
  t.\text{source} = s \;\wedge\; t.\text{target} = \text{Recovered}
\end{equation}

\noindent
This holds for all five designs. In OP\_VAULT, recovery requires the
\texttt{recoveryauth} key---the transition exists but is gated by key possession.

\paragraph{Summary of results.}
The Alloy Analyzer checked 41 assertions and found satisfying instances for
25 scenario predicates. Table~\ref{tab:alloy_results} summarizes the results.

\begin{table}[t]
\centering
\caption{Alloy verification results. 34 of 35 safety checks pass; the one
failure confirms cold key extraction in CAT+CSFS (TM11).}
\label{tab:alloy_results}
\small
\begin{tabular}{@{}lccp{4.5cm}@{}}
\toprule
Model & Checks & Scenarios & Key findings \\
\midrule
CTV        & 6 pass & 3 instances &
  Address reuse produces stuck UTXO \\
CCV        & 7 pass & 5 instances &
  Griefing reachable; mode confusion contained \\
OP\_VAULT  & 8 pass & 6 instances &
  Dual-key = DoS not theft; key-loss = permanent \\
CAT+CSFS   & 5 pass, 1 fail$^*$ & 3 instances &
  Cold key extraction (TM11) confirmed \\
Simplicity & 6 pass & 3 instances &
  Recovery constrained; hot key = grief only \\
Cross       & 2 pass & 5 instances &
  No cross-vault injection \\
\midrule
\textbf{Total} & \textbf{34 pass, 1 fail} & \textbf{25 instances} & \\
\bottomrule
\end{tabular}

\vspace{4pt}
{\footnotesize $^*$Expected failure: Alloy finds a counterexample confirming that
cold key compromise enables fund extraction (TM11).}
\end{table}

% ======================================================================
\section{Three-Layer Evidence Methodology} \label{se:three_layer}

The formal verification complements the empirical measurements and economic projections
to form a three-layer evidence methodology (Figure~\ref{fig:three_layer}). Each layer
answers a different question about the same attack:

\begin{figure}[t]
\centering
\begin{tikzpicture}[
  layer/.style={rectangle, draw, rounded corners=4pt, minimum width=10cm,
                minimum height=1.2cm, font=\small, align=center},
  >=stealth, thick,
]
  \node[layer, fill=purple!10] (alloy)
    {\textbf{Layer 1: Alloy (structural)}\\
     ``Is this attack \emph{possible} in the state machine?''};
  \node[layer, fill=green!10, below=0.5cm of alloy] (regtest)
    {\textbf{Layer 2: Regtest (empirical)}\\
     ``\emph{How much} does this attack cost in vsize?''};
  \node[layer, fill=orange!10, below=0.5cm of regtest] (fee)
    {\textbf{Layer 3: Fee model (economic)}\\
     ``\emph{When} is this attack rational at fee rate~$r$?''};

  \draw[->] (alloy) -- node[right, font=\footnotesize] {possible $\rightarrow$ measure} (regtest);
  \draw[->] (regtest) -- node[right, font=\footnotesize] {vsize $\rightarrow$ project} (fee);
\end{tikzpicture}
\caption{Three-layer evidence methodology. Alloy determines structural possibility.
Regtest measures concrete costs. The fee model projects costs into real-world
fee environments to determine economic rationality.}
\label{fig:three_layer}
\end{figure}

\begin{itemize}
\item \textbf{Layer 1 (Alloy): Is the attack structurally possible?}
  If the Alloy check passes (no counterexample), the attack is structurally
  impossible within the modeled scope. If a counterexample is found, the attack
  is possible and proceeds to Layer~2.

\item \textbf{Layer 2 (Regtest): How much does it cost?}
  The regtest experiments measure the concrete vsize of the attack and defense
  transactions, converting a structural possibility into a quantitative cost.

\item \textbf{Layer 3 (Fee model): When is it rational?}
  The fee projection model (Equation~\ref{eq:fee_projection}) converts vsize into
  sats at each fee rate, and the rationality condition
  (Equation~\ref{eq:rationality}) determines when the attack is economically justified.
\end{itemize}

Table~\ref{tab:three_layer} illustrates the alignment for three representative attacks.

\begin{table}[t]
\centering
\caption{Three-layer evidence alignment for representative attacks.}
\label{tab:three_layer}
\small
\begin{tabular}{@{}lp{2.5cm}p{3cm}p{3.5cm}@{}}
\toprule
Attack & Alloy (Layer 1) & Regtest (Layer 2) & Fee model (Layer 3) \\
\midrule
CCV griefing (TM2)
  & Reachable (keyless recovery exists)
  & 122\vB{}/round attacker, 454\vB{}/round defender
  & Rational at all fee rates; $\gamma = 0.27$ \\[6pt]

CTV fee pinning (TM1)
  & Reachable (anchor output spendable)
  & 14{,}300~sats total
  & Rational whenever vault $> 20{,}000$~sats \\[6pt]

CAT+CSFS cold theft (TM11)
  & Counterexample found (cold key extracts)
  & 125\vB{} recovery tx
  & Rational at any fee rate (immediate theft) \\
\bottomrule
\end{tabular}
\end{table}

The three layers reinforce each other: Alloy gives confidence that the regtest
experiments measure the right attacks (no missed paths), the regtest measurements
give concrete costs to the Alloy-identified possibilities, and the fee model
transforms those costs into deployment guidance.
